\chapter{并发编程基础}
\label{ch:12-outline}

并发会同时改变执行顺序和状态交接方式。先确定数据归属、结果传递与退出路径，再选择线程、进程或协程。


\begin{figure}[htbp]
  \centering
  \begin{tikzpicture}[
    book mindmap,
    level 1 concept/.append style={minimum size=22mm,text width=21mm}
  ]
    \node[text=white,minimum size=25mm,text width=19mm] {第 12 章}
      child[grow=163,level distance=38mm,concept color=DeepBlue] {
        node[text=white] {线程与锁}
      }
      child[grow=17,level distance=38mm,concept color=PythonYellow] {
        node {进程与队列}
      }
      child[grow=-90,level distance=28mm,concept color=PythonBlue!22!white] {
        node {异步任务}
      };
  \end{tikzpicture}
  \caption*{线程共享状态、进程隔离数据、协程协作等待。}
\end{figure}


\section{多线程与 GIL}

线程像同一间工作室里的几位工人，能共享桌上的数据，也容易同时伸手改同一份记录。
用锁把一组必须连着完成的状态修改围起来，再用队列或返回值交接结果。

启用 GIL 的 CPython 同一时刻由一个线程执行受它保护的 Python 代码，
等待 I/O 时线程仍有用处。自由线程构建和会释放 GIL 的扩展有不同表现，
判断性能前，先核对解释器与任务类型。GIL 也不能代替业务上的共享状态保护。三种并发方式的交接机制见图~\ref{fig:visual-ch12}。


\begin{lstlisting}[title={\code{locked_count.py}}]
from threading import Lock, Thread

total = 0
lock = Lock()


def count():
    global total
    for _ in range(1000):
        # 把读、加、写作为一段受保护的操作。
        with lock:
            total += 1


workers = [Thread(target=count) for _ in range(2)]
for worker in workers:
    worker.start()
for worker in workers:
    worker.join()
print(total)
\end{lstlisting}


\textbf{运行核对}\enspace 预期输出：
\begin{lstlisting}[language={},numbers=none]
2000
\end{lstlisting}

\par


\section{多进程基础}

进程像另外开一间工作室。各自有内存，任务和结果要经过可传递的边界。
将工作函数放在模块顶层，传入简单数据，再把结果交回主进程。
启动新进程会有成本，先把任务分到合适大小，再比较并行是否划算。

Windows 常用 spawn 启动方式，会重新导入入口模块。
创建进程池的动作放进入口守卫，
在 PyCharm 里也从这份脚本入口运行。
下面用同一批数字核对进程池返回结果的顺序。


\begin{lstlisting}[title={\code{process_squares.py}}]
from concurrent.futures import ProcessPoolExecutor
from multiprocessing import get_context


def square(value):
    # 顶层函数便于子进程导入。
    return value * value


if __name__ == "__main__":
    with ProcessPoolExecutor(
        max_workers=2, mp_context=get_context("spawn")
    ) as pool:
        print(list(pool.map(square, [2, 3, 4])))
\end{lstlisting}


\textbf{运行核对}\enspace 预期输出：
\begin{lstlisting}[language={},numbers=none]
[4, 9, 16]
\end{lstlisting}

\par


\section[异步 I/O]{异步 I/O（asyncio 极简入门）}

异步像点菜后拿着号码去做别的事。事件循环安排协程推进，
协程碰到可等待的操作时用 await 让出执行权。
这要求所等的操作也配合异步接口，直接在协程里做耗时的同步调用仍会卡住循环。

\code{gather()} 可以一起等多个协程。失败、取消和超时都要在调用边界想清楚，
清理放进 finally 或异步上下文管理器。下面先只观察完成顺序和结果顺序，
接着再用本地网络练习处理状态。

\subsection*{动手｜对照完成顺序与结果顺序}

异步任务类似按单号取餐：完成顺序可能不同于提交顺序，但可以按原任务顺序整理返回结果。这里用短暂的异步等待模拟两份检查，便于专心看任务顺序。

\begin{lstlisting}[title={\code{async_checks.py}}]
import asyncio


async def check(name, delay):
    # 等待时让出事件循环，其他任务可以继续。
    await asyncio.sleep(delay)
    print("完成", name)
    return name


async def main():
    # gather 按传入顺序返回各自结果。
    results = await asyncio.gather(
        check("slow", 0.05),
        check("fast", 0.0),
    )
    print(results)


if __name__ == "__main__":
    asyncio.run(main())
\end{lstlisting}

\textbf{运行核对}\enspace fast 应先完成，gather 返回的列表仍按传入顺序排列。把其中的异步等待换成阻塞睡眠，再对照这两种顺序。

\begin{myunderstanding}
线程共享内存，进程通过传递数据协作，协程在等待处交还执行权。先看任务在等什么，再选并发方式。
\end{myunderstanding}

\begin{pycharmpractice}
多进程脚本从带入口守卫的文件运行。调试线程时先看当前选中的线程，再检查共享变量和锁的位置。
\end{pycharmpractice}

\begin{pitfall}
协程里直接调用同步阻塞函数会拖住事件循环。异步代码也要检查所调用函数是否真正提供异步接口。
\end{pitfall}
协程暂停与恢复的顺序见图~\ref{fig:chapter-12-flow}，等待时仍有其他任务可以推进。
\begin{figure}[htbp]
  \centering
  \begin{tikzpicture}[node distance=5mm]
    \node[python flow step,text width=25mm] (a) {提交协程};
    \node[python flow step,text width=25mm,right=of a] (b) {运行到 await};
    \node[python flow step,text width=25mm,right=of b] (c) {等待外部完成};
    \node[python flow step,text width=25mm,right=of c] (d) {恢复协程};
    \draw[python flow arrow] (a) -- (b);
    \draw[python flow arrow] (b) -- (c);
    \draw[python flow arrow] (c) -- (d);
  \end{tikzpicture}
  \caption{协程在异步等待处让出执行权，事件循环可以推进其他就绪任务。}
  \label{fig:chapter-12-flow}
\end{figure}

\subsection*{进阶边界演练}
取消任务也要走完清理路径。协程里的 finally 适合归还已经取得的资源。一次任务失败以后，其他任务是否仍继续，要按所用的 gather、TaskGroup 或线程池接口来核对，不能只凭相似名字推断。

\begin{figure}[htbp]
\centering
\BookImage{concurrency}{\linewidth}
\caption{线程、进程和协程的交接机制示意，条带长度不表示实测耗时。}
\label{fig:visual-ch12}
\end{figure}

\section{异步任务的生命周期：启动、超时与取消}\label{rev:python:async}
调用 async 函数只会生成协程对象。await 表示等待可等待对象，create\_task 才把协程安排为任务。事件循环在任务主动让出执行权时调度其他任务。

脚本入口使用 \code{asyncio.run(main())}。它创建事件循环，并在退出时关闭。已有事件循环的笔记本环境应在顶层使用 await，不能嵌套调用 run。

\subsection{结构化并发与任务收尾}
TaskGroup 的上下文退出前会等待它创建的任务。某任务以 CancelledError 以外的异常失败时，组会取消其余任务并在收尾后汇总异常；需要按组处理时在外层使用 except*。默认 gather 的第一个异常向调用者传播时，其他任务不保证被取消，不能把两个接口当成完全相同的清理语义。
\code{asyncio.timeout} 限制上下文耗时，在超时边界外捕获 TimeoutError。任务取消通过 CancelledError 在可中断点传递；资源释放应放进 finally 或 async with。一般不吞掉 CancelledError，若显式调用 cancel，应继续 await 任务收集退出结果。
\lstinputlisting[title={\code{examples/revision/async_lifecycle.py}}]{examples/revision/async_lifecycle.py}
预期输出 fast 为 ok、slow 为 timeout，最后显示取消已收集。这里用 sleep 模拟等待，不请求外部服务。真实并发还要用 Semaphore 限制在途任务，并为连接池、队列长度和响应体大小设置边界。

\begin{pitfall}
在协程里调用阻塞的 requests 或 time.sleep 会阻塞事件循环。可改用异步客户端，或对适合在线程执行的阻塞 I/O 使用 asyncio.to\_thread。CPU 密集任务应另行评估进程或原生扩展，并发任务数不是越大越好。
\end{pitfall}
API 与取消语义以 \href{https://docs.python.org/3.11/library/asyncio-task.html}{Python asyncio 任务文档} 为准。

\section{自由线程构建：版本、扩展与共享状态}
Python 3.13 首次提供实验性的自由线程构建；Python 3.14 将该构建列为正式支持，仍是可选构建，不能据此说普通 CPython 已默认去掉 GIL。使用时检查解释器的构建信息和扩展兼容性；不支持自由线程的扩展可能重新启用 GIL。
自由线程不等于业务操作自动具有原子性：例如“先判断键是否存在，再写入”的组合仍需锁或其他同步机制。线程是否加速取决于任务、扩展、共享数据竞争与运行环境，应以实际基准测试判断。
版本依据见 \href{https://docs.python.org/3.14/whatsnew/3.14.html}{Python 3.14 官方版本说明}。

